\chapter*{Preface}
\addcontentsline{toc}{chapter}{Preface}

\section{Purpose of This Manual}

This manual serves as the comprehensive reference for the OmniLaTeX document
class---a modular, extensible \LaTeX{} template designed to unify the production
of academic and technical documents under a single framework.  Rather than
maintaining separate templates for theses, articles, reports, presentations, and
posters, OmniLaTeX provides a single class with a rich doctype system that
adapts its behaviour to the task at hand.

The manual documents every public interface offered by the class: class
options, document-type profiles, custom commands, supported packages, and
recommended workflows.  It is intended to be both a tutorial for newcomers and
a definitive reference for experienced users.

\section{Target Audience}

This manual is written for anyone who produces technical or scientific documents
with \LaTeX{}:

\begin{itemize}
    \item \textbf{Graduate students} writing theses or dissertations that must
          conform to institutional formatting requirements.
    \item \textbf{Researchers and academics} preparing journal articles,
          conference papers, and technical reports.
    \item \textbf{Engineers} generating design documents, specifications, and
          project documentation.
    \item \textbf{Technical writers} who need a consistent, maintainable
          template for multi-chapter publications.
\end{itemize}

Readers are assumed to have basic familiarity with \LaTeX{} (e.g.\ they know
what a preamble is and how to compile a document).  No prior experience with
KOMA-Script or OmniLaTeX is required.

\section{How to Use This Book}

The manual is divided into eleven parts, each covering a major topic area.
Chapters are designed to be as self-contained as possible so that readers can
jump directly to the section relevant to their current task.  Where chapters
depend on material presented elsewhere, explicit cross-references are provided.

A recommended reading order for first-time users:

\begin{enumerate}
    \item Part~I (\emph{Foundations}) --- installation, quick start, build
          system, and class options.
    \item The chapter corresponding to your document type in Part~II
          (\emph{Document Types}).
    \item Remaining parts as needed for typography, layout, figures, references,
          and advanced features.
\end{enumerate}

\section{Scope and Relation to the TUHH Thesis Template}

OmniLaTeX originated from the Hamburg University of Technology (TUHH) thesis
template but has since evolved into a general-purpose framework.  This manual
covers the full OmniLaTeX class in its current form.  Features that are
specific to the TUHH institutional profile are noted where relevant but are not
the primary focus.

If you are writing a TUHH thesis, consult the \texttt{doctype=thesis} profile
documentation in Chapter~\ref{ch:theses} for institution-specific configuration
options.

\section{Acknowledgments}

OmniLaTeX builds upon the work of countless contributors to the \TeX{} and
\LaTeX{} ecosystems.  Special thanks go to the KOMA-Script team for the
\texttt{scrbook}, \texttt{scrreprt}, and \texttt{scrartcl} base classes that
form the foundation of every OmniLaTeX document.  The \texttt{minted},
\texttt{biblatex}, \texttt{glossaries-extra}, and \texttt{tikz} communities
have also been instrumental in making this template possible.

\section{Conventions Used in This Manual}

The following typographic conventions are used throughout:

\begin{itemize}
    \item \texttt{Monospace text} denotes \LaTeX{} commands, environments,
          file names, and code snippets.
    \item \textsf{Sans-serif text} denotes option names and key-value pairs.
    \item \textbf{Bold text} highlights key terms on their first appearance.
    \item \fbox{\small UI paths} indicates clickable menu or interface elements.
\end{itemize}

All code examples are shown in monospaced listings and are designed to be
directly compilable.  Where an example is abbreviated for clarity, the
omissions are explicitly noted.

Full details on these conventions are given in the next chapter,
\emph{Notation \& Conventions}.
